\section{Model Performance Comparison}\label{sec:model-performance-comparison}
For a statistical comparison of the performance of the \gls{cnn} and ensemble across patients, each patient was treated as a paired observation, and it was evaluated whether significant outperformance was achieved by one model over the other.

Because both models were tested on the same set of patients, their performance metrics (e.g., ROC AUC or \gls{eb} score) were not independent but paired.
Additionally, a normal distribution of performance metric across patients was not guaranteed, and the sample size was limited (n~=~12 patients).
Therefore, a non-parametric paired test, namely the Wilcoxon signed-rank test, was used.
Unlike parametric alternatives (e.g., paired t-test), greater robustness to outliers and skewed distributions was offered by this test, making it a safer and more reliable choice in this setting.

By applying this test separately to the relative \gls{tifw}, \gls{eb} sensitivity, \gls{eb} score, and \gls{roc auc}, it was assessed whether the median difference in performance between the two models deviated significantly from zero, thereby providing evidence for systematic superiority of one model over the other.

Now, with trained and evaluated models, the two pillars of evidence for seizure propensity are detailed in the following sections.
